\section{Item-level and round-level mechanisms}
\label{sec:items}

The preceding sections operate at cell level. This section reports the item-level and
agent-round-level evidence that bears on why the cell-level effects occur.

\subsection{Item-difficulty structure}

Table~\ref{tab:item-difficulty} characterises the difficulty spectrum of each benchmark
over the 858{,}089 item rows.

\intable{tab_item_difficulty}

Items that no configuration solves and items that every configuration solves carry no
within-question information and are dropped by the item-level fixed-effect models. The
share of each is a measure of how much of the benchmark is actually discriminating between
the systems under test, and it differs substantially by suite.

\subsection{Confident and wrong: the debate failure mode}

The mechanism proposed by \citet{wynn2025debate} is that communication moves agents toward
agreement on incorrect answers. The direct test is the probability that a team is unanimous
and wrong, contrasted between decentralized debate and independent voting on matched cells.

The pooled gap is $+0.0347$ $[0.0259, 0.0432]$ in the 0.6B and 1.7B band, $+0.0172$
$[0.0108, 0.0235]$ in the 4B and 8B band, and $+0.0164$ $[0.0112, 0.0212]$ in the 14B and
32B band. All three intervals exclude zero and all three are positive: debate produces more
unanimous incorrect consensus than independent voting at every capacity.

The effect more than halves from the smallest band to the largest. Debate-induced
consensus on wrong answers is worst where the models are least able to evaluate a peer's
argument, and diminishes but does not disappear as capacity grows.

This coexists with decentralized having the best average accuracy and the lowest error
amplification of the three multi-agent topologies. The two statistics count different
events: error amplification measures the total error rate, and unanimous-and-wrong measures
a specific mode of failure in which the system is both wrong and maximally confident. Debate
reduces the first while increasing the second, which is the worst possible combination for a
system that exposes agreement as confidence, and it is exactly what the vote-delta results
of Section~\ref{sec:calibration} report at cell level.

The contrast has a structural confound: decentralized runs two rounds and independent runs
one, so the comparison includes an additional round of generation compute. A one-round
decentralized condition would separate communication from compute, and is recorded as a
follow-up.

\subsection{Round dynamics}

Table~\ref{tab:round-accuracy} gives per-agent accuracy by topology and round pooled over
the sweep, and Table~\ref{tab:round-accuracy-size} resolves the two communicating
topologies by capacity. Figure~\ref{fig:round-dynamics-size} plots the same.

\intable{tab_round_accuracy}

\begin{figure}[H]
\centering
\includegraphics[width=0.82\linewidth,height=0.28\textheight,keepaspectratio]{fig_round_dynamics_size.pdf}
\caption{Per-agent accuracy in round 0 and round 1 for the two communicating topologies,
resolved by capacity.}
\label{fig:round-dynamics-size}
\end{figure}

Individual agents improve between rounds. Decentralized agents go from 0.5514 to 0.5741, a
gain of 0.0228, and centralized sub-agents go from 0.5538 to 0.5721, a gain of 0.0184, both
over roughly 835{,}000 to 868{,}000 agent-round observations. Single-agent and independent
cells run one round and provide the reference points at 0.5455 and 0.5553.

Round-zero accuracy is essentially identical across all four topologies, between 0.5455 and
0.5553, which is the expected result since round zero has no peer context and the agents
differ only in their sampling seeds. The between-round gain is therefore attributable to the
peer information rather than to a difference in the agents themselves.

\intable{tab_round_accuracy_size}

The capacity-resolved version identifies where the gain comes from and is not available in
the existing figure set. It is the natural companion to the confident-wrong result: the same
communication that raises individual round-one accuracy also raises the rate of unanimous
incorrect consensus, and both effects are strongest at small capacity.

\subsection{Centralized anatomy}

Within centralized cells the orchestrator and the sub-agents can be compared directly. The
orchestrator's own logprob confidence beats the vote fraction over its sub-agents in 4 of 6
strata, a weak and unreliable advantage. This is the cell-level basis for the conclusion in
Section~\ref{sec:contrasts} that centralized orchestration is not a default improvement: it
neither delivers better accuracy than debate nor a reliably better confidence signal than
voting, while costing the same six generation turns as debate.

\subsection{Context sharing at item level}

The item-level test of context sharing confirms the cell-level split. The accuracy component
is regime-dependent, with the size interaction taking the opposite sign to the one
hypothesised, at $-0.0100$. The convergence component is supported: exposing chain of
thought tightens the spread of confidence across agents in five of six topology and
capacity-band strata.

Tightening confidence spread without improving individual calibration is the item-level
signature of the cell-level result that plus-CoT raises vote calibration error in 94 of 119
controlled strata. Agents converge on a common view; that convergence is not accompanied by
a corresponding gain in the reliability of the view.

\subsection{A corrected estimate}

The item-level final-producer delta reported here differs from an earlier internal draft by
roughly a factor of eight. The earlier version used the mean per-agent confidence as the
baseline in Equation~\ref{eq:estimand} instead of the per-agent calibration error, which
inflated the estimate; the GPQA centralized value moved from $+0.059$ to $+0.004$ once
corrected. The corrected values are the ones reported throughout this document, and they are
consistent with the cell-level range of $+0.004$ to $+0.020$.

The direction of the correction matters for the conclusion. The uncorrected numbers would
have supported a claim that coordination substantially degrades the decisive model's
calibration. The corrected numbers do not.

\subsection{Interactions and configuration ranking}

Two-way interaction models, gradient-boosted importance estimates with grouped
cross-validation, and Friedman interaction statistics were fitted over the pooled core grid;
the resulting figures are in Appendix~\ref{app:atlas}. Capacity and reasoning dominate the
importance rankings, consistent with the marginal effect sizes.

Configuration-level ranking was also attempted. Cross-benchmark rank correlation of
configuration lift is 0.554 among the knowledge multiple-choice pairs and 0.454 between
those and the mathematical benchmark, and five configurations reach the top decile on at
least three benchmarks.

These rankings are not reliable at present coverage and are reported for completeness only.
The top-ten candidate pool has full three-seed coverage for only 20 percent of its entries
and one-seed coverage for 33 percent, and split-half analysis shows the expected shrinkage
of a top-$k$ selection made under that much replication noise. On GPQA specifically, no
multi-agent configuration shows a positive lift over the best same-capacity single-agent
configuration that survives false-discovery-rate control. No configuration-level
recommendation should be drawn from this sweep until the remaining cells complete.
